\section{RELATIONNEL}

\subsection{Traduction du diagramme Entité-Association}

On a traduit notre diagramme Entité-Association en un schéma relationnel composé de 25 tables. On présente ci-dessous la traduction littérale de chaque table avec ses attributs, clés primaires, clés étrangères et contraintes.

\subsubsection{Entités simples}

\begin{itemize}
    \item \textbf{Contact}(\underline{idContact}, nom, prénom, numTel, email)
    \begin{itemize}
        \item Contraintes : \texttt{nom NOT NULL}, \texttt{prenom NOT NULL}
        \item \textbf{Contraintes non implantables} :
        \begin{itemize}
            \item Un contact est soit celui d'un producteur, soit celui d'un client (exclusivité mutuelle).
        \end{itemize}
    \end{itemize}
    
    \item \textbf{Client}(\underline{idClient}, idContact)
    \begin{itemize}
        \item Clé étrangère : \texttt{idContact} référence \texttt{Contact(idContact)} ON DELETE CASCADE
        \item Contraintes : \texttt{idContact NOT NULL UNIQUE}
        \item \textbf{Contraintes non implantables} :
        \begin{itemize}
            \item Droit à l'oubli : gestion particulière pour supprimer les données tout en conservant l'historique des commandes si nécessaire.
        \end{itemize}
    \end{itemize}
    
    \item \textbf{Adresse}(\underline{idAdresse}, rue, ville, codePostal)
    \begin{itemize}
        \item Contraintes : \texttt{rue NOT NULL}, \texttt{ville NOT NULL}, \texttt{codePostal NOT NULL}
        \item \textbf{Contraintes non implantables} :
        \begin{itemize}
            \item Une adresse est soit celle d'un producteur, soit celle d'un client (via PossedeAdresse).
        \end{itemize}
    \end{itemize}
    
    \item \textbf{Producteur}(\underline{idProducteur}, lattitude, longitude, idAdresse, idContact)
    \begin{itemize}
        \item Clé étrangère : \texttt{idAdresse} référence \texttt{Adresse(idAdresse)}
        \item Clé étrangère : \texttt{idContact} référence \texttt{Contact(idContact)} ON DELETE CASCADE
        \item Contraintes : \texttt{idAdresse NOT NULL}, \texttt{idContact NOT NULL UNIQUE}
    \end{itemize}
    
    \item \textbf{TypeActivite}(\underline{typeActivite})
    
    \item \textbf{Contenant}(\underline{referenceContenant}, typeContenant, capaciteContenant, stockContenant, caractereContenant, prixContenant)
    \begin{itemize}
        \item Contraintes : \texttt{typeContenant NOT NULL}, \texttt{capaciteContenant NOT NULL}, \texttt{stockContenant NOT NULL}, \texttt{caractereContenant NOT NULL}, \texttt{prixContenant NOT NULL}
        \item Contraintes textuelles : \texttt{capaciteContenant > 0}, \texttt{stockContenant >= 0}, \texttt{caractereContenant IN ('Réutilisable', 'Non-Réutilisable')}, \texttt{prixContenant >= 0}
    \end{itemize}
    
    \item \textbf{Disponibilite}(\underline{idDisponibilite}, debutDisponibilite, finDisponibilite, statutProduit)
    \begin{itemize}
        \item Contraintes : \texttt{debutDisponibilite NOT NULL}, \texttt{finDisponibilite NOT NULL}, \texttt{statutProduit NOT NULL}
        \item Contraintes textuelles : \texttt{statutProduit IN ('Disponible', 'Pas disponible')}, \texttt{finDisponibilite >= debutDisponibilite}
    \end{itemize}
    
    \item \textbf{Perte}(\underline{idPerte}, datePerte, naturePerte)
    \begin{itemize}
        \item Contraintes : \texttt{datePerte NOT NULL}, \texttt{naturePerte NOT NULL}
        \item Contraintes textuelles : \texttt{naturePerte IN ('Vol', 'Casse')}
    \end{itemize}
    
    \item \textbf{Produit}(\underline{idProduit}, \underline{idProducteur}, categorie, description, bio, label, nomProduit, allergene, origineGeographique, delaiDisponibilite)
    \begin{itemize}
        \item Clé étrangère : \texttt{idProducteur} référence \texttt{Producteur(idProducteur)} ON DELETE CASCADE
        \item Contraintes : \texttt{categorie NOT NULL}, \texttt{nomProduit NOT NULL}
    \end{itemize}
    
    \item \textbf{Stock}(\underline{idStock}, idProduit, idProducteur)
    \begin{itemize}
        \item Clé étrangère : \texttt{(idProduit, idProducteur)} référence \texttt{Produit(idProduit, idProducteur)} ON DELETE CASCADE
        \item Contraintes : \texttt{idProduit NOT NULL}, \texttt{idProducteur NOT NULL}
    \end{itemize}
    
    \item \textbf{Commande}(\underline{idCommande}, dateCommande, dateRecuperation, datePaiement, heureCommande, statutCommande, modePaiement, modeRecuperation, idClient)
    \begin{itemize}
        \item Clé étrangère : \texttt{idClient} référence \texttt{Client(idClient)} ON DELETE CASCADE
        \item Contraintes : \texttt{dateCommande NOT NULL}, \texttt{heureCommande NOT NULL}, \texttt{statutCommande NOT NULL}, \texttt{modePaiement NOT NULL}, \texttt{modeRecuperation NOT NULL}, \texttt{idClient NOT NULL}
        \item Contraintes textuelles : \texttt{statutCommande IN ('En préparation', 'Prête', 'En livraison', 'Annulée', 'Récupérée/Livrée')}, \texttt{modePaiement IN ('En ligne', 'En Boutique')}, \texttt{modeRecuperation IN ('Boutique', 'Domicile')}
        \item \textbf{Contraintes non implantables} :
        \begin{itemize}
            \item Annulation par le client autorisée tant que la commande est "En préparation".
            \item Paiement en boutique uniquement si le statut est "Prête".
            \item Pas de statut "En préparation" pour une commande avec le mode de récupération "Boutique".
            \item \texttt{dateRecuperation} et \texttt{datePaiement} sont non NULL si \texttt{statutCommande = 'Récupérée/Livrée'}.
            \item Pas de possibilité de payer en boutique si le mode de récupération est "Domicile".
            \item \texttt{dateRecuperation >= dateCommande}.
            \item \texttt{datePaiement <= dateRecuperation} (si les deux sont non NULL).
        \end{itemize}
    \end{itemize}
    
    \item \textbf{ModeRecuperationDomicile}(\underline{idModeRecuperationDomicile}, paysLivraison, poidsTotalCommande, distanceAdresseBoutique, dateEstimeeLivraison, typePaysLivraison, idCommande, idAdresse)
    \begin{itemize}
        \item Clé étrangère : \texttt{idCommande} référence \texttt{Commande(idCommande)} ON DELETE CASCADE
        \item Clé étrangère : \texttt{idAdresse} référence \texttt{Adresse(idAdresse)}
        \item Contraintes : \texttt{paysLivraison NOT NULL}, \texttt{poidsTotalCommande NOT NULL}, \texttt{distanceAdresseBoutique NOT NULL}, \texttt{typePaysLivraison NOT NULL}, \texttt{idCommande NOT NULL UNIQUE}, \texttt{idAdresse NOT NULL}
        \item Contraintes textuelles : \texttt{poidsTotalCommande > 0}, \texttt{distanceAdresseBoutique >= 0}, \texttt{typePaysLivraison IN ('France Métropolitaine', 'DOM-TOM', 'International')}
        \item \textbf{Contraintes non implantables} :
        \begin{itemize}
            \item La date de livraison (\texttt{dateEstimeeLivraison}) doit être communiquée au client.
            \item \texttt{dateEstimeeLivraison >= dateCommande} (de la commande associée).
        \end{itemize}
    \end{itemize}
    
    \item \textbf{LigneCommande}(\underline{idLigneCommande}, \underline{idCommande}, prixUnitaire, sousTotalLigne)
    \begin{itemize}
        \item Clé étrangère : \texttt{idCommande} référence \texttt{Commande(idCommande)} ON DELETE CASCADE
        \item Contraintes : \texttt{prixUnitaire NOT NULL}, \texttt{sousTotalLigne NOT NULL}, \texttt{idCommande NOT NULL}
        \item Contraintes textuelles : \texttt{prixUnitaire >= 0}, \texttt{sousTotalLigne >= 0}
        \item \textbf{Contraintes non implantables} :
        \begin{itemize}
            \item Le total de la commande doit être égal à la somme des \texttt{sousTotalLigne} de toutes les lignes.
        \end{itemize}
    \end{itemize}
\end{itemize}

\subsubsection{Sous-types d'entités}

\begin{itemize}
    \item \textbf{Lot}(\underline{idLot}, dateReception, dateLimite, typeDateLimite, idStock)
    \begin{itemize}
        \item Clé étrangère : \texttt{idStock} référence \texttt{Stock(idStock)} ON DELETE CASCADE
        \item Contraintes : \texttt{dateReception NOT NULL}, \texttt{dateLimite NOT NULL}, \texttt{typeDateLimite NOT NULL}
        \item Contraintes textuelles : \texttt{typeDateLimite IN ('DLC', 'DLUO')}
        \item \textbf{Contraintes non implantables} :
        \begin{itemize}
            \item \texttt{dateLimite >= dateReception}.
            \item Un lot correspond à une livraison à l'épicerie d'un jour particulier.
        \end{itemize}
    \end{itemize}
    
    \item \textbf{LotVrac}(\underline{idLot}, quantiteDisponibleVrac)
    \begin{itemize}
        \item Clé étrangère : \texttt{idLot} référence \texttt{Lot(idLot)} ON DELETE CASCADE
        \item Contraintes : \texttt{quantiteDisponibleVrac NOT NULL}
        \item Contraintes textuelles : \texttt{quantiteDisponibleVrac >= 0}
        \item \textbf{Contraintes non implantables} :
        \begin{itemize}
            \item Un lot ne peut être à la fois vrac ET préconditionné (exclusivité mutuelle).
            \item La quantité disponible doit être cohérente avec les quantités commandées et les pertes.
        \end{itemize}
    \end{itemize}
    
    \item \textbf{LotPreconditionne}(\underline{idLot}, quantiteDisponiblePreconditionne)
    \begin{itemize}
        \item Clé étrangère : \texttt{idLot} référence \texttt{Lot(idLot)} ON DELETE CASCADE
        \item Contraintes : \texttt{quantiteDisponiblePreconditionne NOT NULL}
        \item Contraintes textuelles : \texttt{quantiteDisponiblePreconditionne >= 0}
        \item \textbf{Contraintes non implantables} :
        \begin{itemize}
            \item Un lot ne peut être à la fois vrac ET préconditionné (exclusivité mutuelle).
            \item La quantité disponible doit être cohérente avec les quantités commandées et les pertes.
        \end{itemize}
    \end{itemize}
    
    \item \textbf{Conditionnement}(\underline{idConditionnement}, prixAchatProducteur, prixVenteClient, idProduit, idProducteur)
    \begin{itemize}
        \item Clé étrangère : \texttt{(idProduit, idProducteur)} référence \texttt{Produit(idProduit, idProducteur)} ON DELETE CASCADE
        \item Contraintes : \texttt{prixAchatProducteur NOT NULL}, \texttt{prixVenteClient NOT NULL}
        \item Contraintes textuelles : \texttt{prixAchatProducteur >= 0}, \texttt{prixVenteClient >= 0}
    \end{itemize}
    
    \item \textbf{ConditionnementPreconditionne}(\underline{idConditionnement}, poidsSachet)
    \begin{itemize}
        \item Clé étrangère : \texttt{idConditionnement} référence \texttt{Conditionnement(idConditionnement)} ON DELETE CASCADE
        \item Contraintes : \texttt{poidsSachet NOT NULL}
        \item Contraintes textuelles : \texttt{poidsSachet > 0}
        \item \textbf{Contraintes non implantables} :
        \begin{itemize}
            \item Un conditionnement ne peut être à la fois vrac ET préconditionné (exclusivité mutuelle).
        \end{itemize}
    \end{itemize}
    
    \item \textbf{ConditionnementVrac}(\underline{idConditionnement})
    \begin{itemize}
        \item Clé étrangère : \texttt{idConditionnement} référence \texttt{Conditionnement(idConditionnement)} ON DELETE CASCADE
        \item \textbf{Contraintes non implantables} :
        \begin{itemize}
            \item Un conditionnement ne peut être à la fois vrac ET préconditionné (exclusivité mutuelle).
            \item Un seul conditionnement vrac pour un même produit et producteur.
        \end{itemize}
    \end{itemize}
    
    \item \textbf{LigneCommandeContenant}(\underline{idLigneCommande}, \underline{idCommande}, quantiteCommandeContenant)
    \begin{itemize}
        \item Clé étrangère : \texttt{(idLigneCommande, idCommande)} référence \texttt{LigneCommande(idLigneCommande, idCommande)} ON DELETE CASCADE
        \item Contraintes : \texttt{quantiteCommandeContenant NOT NULL}
        \item Contraintes textuelles : \texttt{quantiteCommandeContenant > 0}
        \item \textbf{Contraintes non implantables} :
        \begin{itemize}
            \item Une ligne de commande ne peut être à la fois produit ET contenant (exclusivité mutuelle).
        \end{itemize}
    \end{itemize}
    
    \item \textbf{LigneCommandeProduit}(\underline{idLigneCommande}, \underline{idCommande}, idProduit, idProducteur)
    \begin{itemize}
        \item Clé étrangère : \texttt{(idLigneCommande, idCommande)} référence \texttt{LigneCommande(idLigneCommande, idCommande)} ON DELETE CASCADE
        \item Clé étrangère : \texttt{(idProduit, idProducteur)} référence \texttt{Produit(idProduit, idProducteur)}
        \item Contraintes : \texttt{idProduit NOT NULL}, \texttt{idProducteur NOT NULL}
        \item \textbf{Contraintes non implantables} :
        \begin{itemize}
            \item Une ligne de commande ne peut être à la fois produit ET contenant (exclusivité mutuelle).
        \end{itemize}
    \end{itemize}
    
    \item \textbf{LigneCommandeProduitVrac}(\underline{idLigneCommande}, \underline{idCommande}, quantiteCommandeVrac)
    \begin{itemize}
        \item Clé étrangère : \texttt{(idLigneCommande, idCommande)} référence \texttt{LigneCommandeProduit(idLigneCommande, idCommande)} ON DELETE CASCADE
        \item Contraintes : \texttt{quantiteCommandeVrac NOT NULL}
        \item Contraintes textuelles : \texttt{quantiteCommandeVrac > 0}
        \item \textbf{Contraintes non implantables} :
        \begin{itemize}
            \item Une ligne de commande produit ne peut être à la fois vrac ET préconditionné (exclusivité mutuelle).
            \item La quantité commandée doit être inférieure ou égale à la quantité disponible dans le lot.
        \end{itemize}
    \end{itemize}
    
    \item \textbf{LigneCommandeProduitPreconditionne}(\underline{idLigneCommande}, \underline{idCommande}, quantiteCommandePreconditionne)
    \begin{itemize}
        \item Clé étrangère : \texttt{(idLigneCommande, idCommande)} référence \texttt{LigneCommandeProduit(idLigneCommande, idCommande)} ON DELETE CASCADE
        \item Contraintes : \texttt{quantiteCommandePreconditionne NOT NULL}
        \item Contraintes textuelles : \texttt{quantiteCommandePreconditionne > 0}
        \item \textbf{Contraintes non implantables} :
        \begin{itemize}
            \item Une ligne de commande produit ne peut être à la fois vrac ET préconditionné (exclusivité mutuelle).
            \item La quantité commandée doit être inférieure ou égale à la quantité disponible dans le lot.
        \end{itemize}
    \end{itemize}
    
    \item \textbf{PerteProduit}(\underline{idPerte}, idProduit, idProducteur)
    \begin{itemize}
        \item Clé étrangère : \texttt{idPerte} référence \texttt{Perte(idPerte)} ON DELETE CASCADE
        \item Clé étrangère : \texttt{(idProduit, idProducteur)} référence \texttt{Produit(idProduit, idProducteur)}
        \item Contraintes : \texttt{idProduit NOT NULL}, \texttt{idProducteur NOT NULL}
        \item \textbf{Contraintes non implantables} :
        \begin{itemize}
            \item Une perte ne peut être à la fois produit ET contenant (exclusivité mutuelle).
        \end{itemize}
    \end{itemize}
    
    \item \textbf{PerteProduitPreconditionne}(\underline{idPerte}, quantitePerduePreconditionne)
    \begin{itemize}
        \item Clé étrangère : \texttt{idPerte} référence \texttt{PerteProduit(idPerte)} ON DELETE CASCADE
        \item Contraintes : \texttt{quantitePerduePreconditionne NOT NULL}
        \item Contraintes textuelles : \texttt{quantitePerduePreconditionne > 0}
        \item \textbf{Contraintes non implantables} :
        \begin{itemize}
            \item Une perte produit ne peut être à la fois vrac ET préconditionné (exclusivité mutuelle).
        \end{itemize}
    \end{itemize}
    
    \item \textbf{PerteProduitVrac}(\underline{idPerte}, quantitePerdueVrac)
    \begin{itemize}
        \item Clé étrangère : \texttt{idPerte} référence \texttt{PerteProduit(idPerte)} ON DELETE CASCADE
        \item Contraintes : \texttt{quantitePerdueVrac NOT NULL}
        \item Contraintes textuelles : \texttt{quantitePerdueVrac > 0}
        \item \textbf{Contraintes non implantables} :
        \begin{itemize}
            \item Une perte produit ne peut être à la fois vrac ET préconditionné (exclusivité mutuelle).
        \end{itemize}
    \end{itemize}
    
    \item \textbf{PerteContenant}(\underline{idPerte}, quantitePerdueContenant, referenceContenant)
    \begin{itemize}
        \item Clé étrangère : \texttt{idPerte} référence \texttt{Perte(idPerte)} ON DELETE CASCADE
        \item Clé étrangère : \texttt{referenceContenant} référence \texttt{Contenant(referenceContenant)}
        \item Contraintes : \texttt{quantitePerdueContenant NOT NULL}, \texttt{referenceContenant NOT NULL}
        \item Contraintes textuelles : \texttt{quantitePerdueContenant > 0}
        \item \textbf{Contraintes non implantables} :
        \begin{itemize}
            \item Une perte ne peut être à la fois produit ET contenant (exclusivité mutuelle).
        \end{itemize}
    \end{itemize}
\end{itemize}

\subsubsection{Associations}

\begin{itemize}
    \item \textbf{PossedeAdresse}(\underline{idClient}, \underline{idAdresse})
    \begin{itemize}
        \item Clé étrangère : \texttt{idClient} référence \texttt{Client(idClient)} ON DELETE CASCADE
        \item Clé étrangère : \texttt{idAdresse} référence \texttt{Adresse(idAdresse)} ON DELETE CASCADE
    \end{itemize}
    
    \item \textbf{Exerce}(\underline{idProducteur}, \underline{typeActivite})
    \begin{itemize}
        \item Clé étrangère : \texttt{idProducteur} référence \texttt{Producteur(idProducteur)} ON DELETE CASCADE
        \item Clé étrangère : \texttt{typeActivite} référence \texttt{TypeActivite(typeActivite)} ON DELETE CASCADE
    \end{itemize}
    
    \item \textbf{ProduitEstDisponible}(\underline{idProduit}, \underline{idProducteur}, \underline{idDisponibilite})
    \begin{itemize}
        \item Clé étrangère : \texttt{(idProduit, idProducteur)} référence \texttt{Produit(idProduit, idProducteur)} ON DELETE CASCADE
        \item Clé étrangère : \texttt{idDisponibilite} référence \texttt{Disponibilite(idDisponibilite)} ON DELETE CASCADE
    \end{itemize}
    
    \item \textbf{PeutEtreUtiliseAvec}(\underline{idConditionnement}, \underline{referenceContenant})
    \begin{itemize}
        \item Clé étrangère : \texttt{idConditionnement} référence \texttt{ConditionnementVrac(idConditionnement)} ON DELETE CASCADE
        \item Clé étrangère : \texttt{referenceContenant} référence \texttt{Contenant(referenceContenant)} ON DELETE CASCADE
    \end{itemize}
\end{itemize}

\subsubsection{Justification des choix pour les sous-types d'entités}

On a utilisé la technique de spécialisation par table pour gérer les relations d'héritage. Cette approche consiste à créer une table pour la classe mère et des tables séparées pour chaque sous-classe, avec une clé étrangère de la sous-classe vers la classe mère.

\paragraph{Lot $\rightarrow$ LotVrac / LotPreconditionne}

On a choisi cette modélisation car les attributs spécifiques sont différents (\texttt{quantiteDisponibleVrac} FLOAT vs \texttt{quantiteDisponiblePreconditionne} INTEGER), les types de données diffèrent, et un lot ne peut être qu'en vrac OU préconditionné (exclusivité mutuelle). La table \texttt{Lot} contient les attributs communs, et les tables spécialisées héritent via \texttt{idLot} qui est à la fois leur clé primaire et une clé étrangère vers \texttt{Lot}.

\paragraph{Conditionnement $\rightarrow$ ConditionnementVrac / ConditionnementPreconditionne}

Même principe : \texttt{Conditionnement} contient les attributs communs, \texttt{ConditionnementPreconditionne} ajoute \texttt{poidsSachet}, et \texttt{ConditionnementVrac} permet de distinguer les conditionnements vrac pour les lier avec des contenants. Un conditionnement ne peut être qu'en vrac OU préconditionné (exclusivité mutuelle).

\paragraph{LigneCommande $\rightarrow$ LigneCommandeProduit / LigneCommandeContenant $\rightarrow$ LigneCommandeProduitVrac / LigneCommandeProduitPreconditionne}

On a créé une hiérarchie à deux niveaux : \texttt{LigneCommande} est spécialisée en \texttt{LigneCommandeProduit} (pour les produits) et \texttt{LigneCommandeContenant} (pour les contenants), puis \texttt{LigneCommandeProduit} est spécialisée en \texttt{LigneCommandeProduitVrac} et \texttt{LigneCommandeProduitPreconditionne}. Cette modélisation permet de gérer les différences d'attributs et de traitements selon le type de ligne.

\paragraph{Perte $\rightarrow$ PerteProduit / PerteContenant $\rightarrow$ PerteProduitVrac / PerteProduitPreconditionne}

Même principe appliqué pour distinguer les pertes de produits (qui nécessitent une référence au produit) des pertes de contenants (qui nécessitent une référence au contenant).

\subsection{Formes normales}

\subsubsection{Troisième forme normale (3NF)}

On a vérifié que nos tables respectent la 3NF. Les attributs non-clés dépendent directement de la clé primaire, sans dépendances transitives problématiques. On note que pour \texttt{Adresse}, on a choisi de ne pas modéliser la dépendance \texttt{codePostal} $\rightarrow$ \texttt{ville} car un code postal peut correspondre à plusieurs villes et on veut permettre la flexibilité pour les adresses internationales.

\subsubsection{Forme normale de Boyce-Codd (BCNF)}

La plupart de nos tables sont presque en BCNF. Les cas particuliers sont \texttt{Adresse} (dépendance \texttt{codePostal} $\rightarrow$ \texttt{ville} non modélisée pour flexibilité) et \texttt{ModeRecuperationDomicile} (dépendance \texttt{paysLivraison} $\rightarrow$ \texttt{typePaysLivraison} conservée pour des raisons pratiques).

\subsection{Contraintes applicatives non implantables}

\subsubsection{Contraintes temporelles}

\begin{itemize}
    \item \texttt{dateRecuperation >= dateCommande}
    \item \texttt{datePaiement <= dateRecuperation} (si les deux sont non NULL)
    \item \texttt{dateLimite >= dateReception} (pour Lot)
    \item \texttt{dateEstimeeLivraison >= dateCommande} (pour ModeRecuperationDomicile)
\end{itemize}

\subsubsection{Contraintes métier complexes}

\begin{itemize}
    \item Annulation par le client autorisée tant que la commande est "En préparation"
    \item Paiement en boutique uniquement si le statut est "Prête"
    \item Pas de statut "En préparation" pour une commande avec le mode de récupération "Boutique"
    \item \texttt{dateRecuperation} et \texttt{datePaiement} sont non NULL si \texttt{statutCommande = 'Récupérée/Livrée'}
    \item Pas de possibilité de payer en boutique si le mode de récupération est "Domicile"
    \item Un contact est soit celui d'un producteur, soit celui d'un client (exclusivité mutuelle)
    \item Une adresse est soit celle d'un producteur, soit celle d'un client
    \item Un lot correspond à une livraison à l'épicerie d'un jour particulier
    \item Un seul conditionnement vrac pour un même produit et producteur
    \item La date de livraison doit être communiquée au client
\end{itemize}

\subsubsection{Contraintes de cohérence transactionnelle}

\begin{itemize}
    \item Stock suffisant avant de valider une commande
    \item Cohérence des quantités : \texttt{quantiteCommandeVrac} <= \texttt{quantiteDisponibleVrac} du lot correspondant
    \item Total de la commande = somme des \texttt{sousTotalLigne}
    \item Quantité disponible cohérente avec les quantités commandées et les pertes
\end{itemize}

\subsubsection{Gestion de la disponibilité}

La disponibilité d'un produit dépend de plusieurs facteurs :
\begin{itemize}
    \item Pour les produits en stock : disponibilité immédiate si quantité suffisante, indisponible si hors saison avec indication de la prochaine période
    \item Pour les produits sur commande : délai de disponibilité spécifié par produit (ex: 72h, 1 semaine), client informé du délai
    \item Le nombre de produits disponibles en stock doit être calculé en tenant compte des quantités commandées et des pertes
\end{itemize}

\subsubsection{Contraintes de cardinalité}

\begin{itemize}
    \item Commande non vide (au moins une ligne)
    \item Stock avec au moins un lot
    \item Produit avec au moins un conditionnement
    \item Client avec au moins une adresse
\end{itemize}

\subsubsection{Contraintes d'intégrité référentielle supplémentaires}

\begin{itemize}
    \item \texttt{ModeRecuperationDomicile} uniquement si \texttt{modeRecuperation = "Domicile"}
    \item Conservation de l'historique des commandes passées même si les produits sont supprimés
    \item Gestion des suppressions en cascade selon les contraintes métier (ex: annuler plutôt que supprimer les commandes en cours)
\end{itemize}
